% ---------------------------------------------------------------------
\section{欧氏 OPE 与凝聚}\label{sec:ee-ope}

微扰论不能告诉我们 $D(Q^2)$ 与强子世界的偏差有多大。算符乘积展开（OPE）把偏差组织成 $1/Q^2$ 的幂次：系数由短程物理决定，长程物理进入算符的矩阵元，对流关联函数即真空期望值（“凝聚”）。本节区分微扰论中的 OPE（\tagA）、凝聚的非微扰意义（\tagC）与在复 $q^2$ 处使用 OPE 的额外假设。

\begin{theorem}[流关联函数的欧氏 OPE \tagA]\label{thm:ee-ope}
在重整化微扰论中，对带红外正规化部分子外态的 Green 函数，当欧氏的 $Q^2=-q^2\to\infty$ 时，
\begin{equation}\label{eq:ee-ope}
\begin{split}
\ii\!\int\!\dd^4x\,\ee^{\ii q\cdot x}\,T\{j^\mu(x)j^\nu(0)\}&=(q^\mu q^\nu-q^2g^{\mu\nu})\Big[C_{\mathbb 1}(Q^2)\,\mathbb 1+\sum_qC_q(Q^2)\,\frac{m_q\bar qq}{Q^4}\\
&\qquad+C_G(Q^2)\,\frac{1}{Q^4}\,\frac{\as}\pi G^a_{\rho\sigma}G^{a\rho\sigma}\Big]+\order{Q^{-6}},
\end{split}
\end{equation}
理解为插入 Green 函数后的渐近展开（余项可乘以 $\ln Q^2$ 的幂；量纲 $\le4$ 的 BRST 恰当算符与运动方程算符的物理矩阵元为零，未写出）；系数与外态及红外正规化无关。对无质量夸克展开中没有 $1/Q^2$ 项，因为不存在量纲为 $2$ 的规范不变局域算符；夸克有质量时 $1/Q^2$ 项只来自 $C_{\mathbb 1}\ni-\frac{N_ce_q^2}{2\pi^2}\frac{m_q^2}{Q^2}$。领头阶系数为
\begin{equation}\label{eq:ee-ope-coeff}
C_q=2e_q^2\big[1+\order{\as}\big],\qquad C_G=\frac1{12}\sum_qe_q^2\big[1+\order{\as}\big]=\frac{R_0}{12N_c}\big[1+\order{\as}\big],
\end{equation}
$C_{\mathbb 1}=\Pi_{\rm PT}$ 是微扰的关联函数（\cref{prop:ee-twoloop}），归一化与\cref{eq:ee-Pidef}相同。% VF-CONV-12
（本定理只涉及部分子态的 Green 函数。）
\uses{thm:prelim-composite-renorm,thm:ee-euclid-ir,def:ee-adler,prop:ee-twoloop}
\end{theorem}
\begin{proofsketch}
(1)~\emph{Zimmermann 恒等式与 Weinberg 定理。}对规范不变流的欧氏两点函数，在 BPHZ 或 BPHZL 重整化的微扰论中，编时乘积减去展开到量纲 $N$ 的正规乘积之和，被积函数有改善的紫外幂次；Weinberg 定理逐图给出余项界 $\order{Q^{2-N}}$（取对数）。非例外欧氏点处的红外有限性由\cref{thm:ee-euclid-ir}保证。
(2)~\emph{协变性与量纲。}流守恒给出 $(q^\mu q^\nu-q^2g^{\mu\nu})$；量纲为 $d$ 的算符的系数按 $(Q^2)^{-d/2}$ 乘以 $\ln Q^2$ 的多项式标度。量纲 $\le4$ 的 Lorentz 标量、规范不变、宇称偶的局域算符只有 $\mathbb 1$、$m_q\bar qq$（手征性使 $\bar qq$ 与 $m_q$ 同时出现）与 $G^aG^a$；$A^a_\mu A^{a\mu}$ 不规范不变。
(3)~\emph{规范变算符。}由\cref{thm:prelim-composite-renorm}（\tagR，$R_{\mathrm{PT}}$），规范不变算符在重整化下只与规范不变算符、BRST 恰当算符和运动方程算符混合；后两类算符在物理态之间的矩阵元为零，因此不必在\cref{eq:ee-ope}中写出。（在微扰论中，量纲为正的算符的无质量真空期望值都是无标度积分，本身为零；在非微扰的真空中 BRST 恰当算符的期望值为零，则属于\cref{asm:intro-np-ope}的内容。）
(4)~\emph{领头系数。}$m_q^2$ 项来自有质量单圈 $\Pi=-\frac{N_ce_q^2}{2\pi^2}\int_0^1\dd x\,x(1-x)\ln\frac{m_q^2+x(1-x)Q^2}{\mu^2}$ 的展开；$C_q$ 由把一条夸克线换成非定域凝聚 $\langle\bar q^\beta(0)q^\alpha(x)\rangle=\frac{\langle\bar qq\rangle}{4N_c}[\delta^{\alpha\beta}-\frac{\ii m_q}4\sla x^{\alpha\beta}+\cdots]$（第二项来自运动方程 $\ii\sla\partial q=m_qq$）得到，四个贡献之和自动横向；$C_G$ 由两个软胶子的前向振幅对胶子动量展开到二阶、再取常数场强的真空平均 $\langle g^2G^a_{\alpha\beta}G^b_{\gamma\delta}\rangle=\frac{\delta^{ab}(g_{\alpha\gamma}g_{\beta\delta}-g_{\alpha\delta}g_{\beta\gamma})}{12(N_c^2-1)}\langle g^2G^2\rangle$ 得到。我们独立完成了这两项计算（含 Ward 恒等式与平均方式无关性的检验），与 Shifman--Vainshtein--Zakharov\cite{ShifmanVainshteinZakharov1979a,ShifmanVainshteinZakharov1979b}及\cite[式~(45)]{Colangelo:2000dp}（$\Pi$ 的定义与本书相同）一致。完整论证见\cite{Zimmermann1973}（标量理论）、\cite[第~12.8~节]{Collins1984}与\cref{sec:appM-massless}。
\end{proofsketch}

\paragraph{地位与缺口}OPE 在重整化微扰论中的全阶证明，对标量理论由 Zimmermann 给出\cite{Zimmermann1973}，对纯 Yang--Mills 的规范不变算符由 Hollands 给出\cite{Hollands:2007zg}；对含夸克的 QCD，Collins 断言它“is proved in a gauge theory the same way as in any other theory”\cite[第~12.8~节]{Collins1984}（在规范理论中与在其他理论中以同样的方式证明），Izubuchi 等人称之为广泛采用的假定\cite[第~II~节]{IzubuchiEtAl2018}——因此是全阶论证（\tagA）。最弱的一步是含夸克与鬼场的规范理论中 Zimmermann 恒等式与完整算符基（规范变算符、鬼场算符、evanescent 算符）的处理（逐图的点态界 E-P，仅在所引定理适用之处）；无质量 QCD 中的细节与缺口见\cref{sec:appM-massless}，完整的地位讨论在后文\cref{thm:dis-ope}（\tagA）处重复并引用本处，OPE 的严格结果汇总于\cref{rem:front-ope-qcd-status}。真空 Adler 函数的 OPE 没有 $1/Q^2$ 项；凝聚项只在\cref{asm:intro-np-ope}之下有非微扰意义——在微扰论中无质量真空期望值都是无标度的，为零。提升为 \tagR 需要对含夸克 QCD 的规范不变流写出 Zimmermann 水平的完整证明。

取真空期望值并换成 Adler 函数（$\Pi\ni C/Q^4$ 给出 $D\ni24\pi^2C/Q^4$），得
\begin{equation}\label{eq:ee-adler-ope}
D(Q^2)=R_0\,d\big(a(Q^2)\big)+\frac{1}{Q^4}\sum_qe_q^2\Big[48\pi^2\langle m_q\bar qq\rangle+2\pi^2\Big\langle\frac{\as}\pi G^2\Big\rangle\Big]+\order{Q^{-6}} ;
\end{equation}
对 SU(3)，胶子凝聚项为 $R_0\frac{2\pi^2}3\langle\frac{\as}\pi G^2\rangle/Q^4$。其非微扰内容就是下面的猜想。

\begin{conjecture}[非微扰 OPE 与凝聚 \tagC]\label{conj:ee-np-ope}
\cref{thm:ee-ope}的展开在 QCD 的真空态之间非微扰地成立，系数不变：$\Pi(-Q^2)=C_{\mathbb 1}+\sum_qC_q\frac{\langle m_q\bar qq\rangle}{Q^4}+C_G\frac{\langle\frac{\as}\pi G^2\rangle}{Q^4}+\order{Q^{-6}}$（$Q^2\to\infty$ 的渐近展开），凝聚是非微扰参数。这是\cref{asm:intro-np-ope}的真空部分。$\langle G^2\rangle$ 的定义不能比 $C_{\mathbb 1}$ 的领头红外重整化子（Borel 平面上 $u=2$ 的极点）的歧义更精确。
\uses{thm:ee-ope}
\end{conjecture}

\paragraph{地位与缺口}QCD 中没有 OPE 的非微扰定义（\tagC）。在代数量子场论中 OPE 可以作为相空间性质的推论得到\cite{Bostelmann:2005rr}，但所需的性质在 QCD 中未被建立。Novikov 等人在模型中讨论了 OPE 是否会失效以及用归一化点 $\mu$ 分离短程与长程贡献\cite{Novikov:1984rf}；David 指出凝聚的定义与红外重整化子的歧义纠缠在一起\cite{David:1983gz,David:1985xj}。与之相容的是：大 $\beta_0$ 极限下 Adler 函数的 Borel 变换在 $u=1$ 处没有红外重整化子极点，第一个红外极点在 $u=2$，对应胶子凝聚（\cref{prop:power-adler-large-beta0}，\tagP，只在该极限中成立）。至于 OPE 作为幂次展开本身的收敛性（而不仅是渐近性），严格结果只有欧氏 $\varphi^4_4$ 理论在每个固定圈阶上的收敛性（有质量情形\cite{Hollands:2011gf}，无质量情形\cite{Holland:2014pna}）——这是固定圈阶的陈述，不涉及非微扰的凝聚；格点上的数值结果是证据而不是证明。解除这一猜想需要 QCD 中 OPE 的非微扰定义与证明。

\begin{remark}[复 $q^2$ 处的 OPE 是额外假设]\label{rem:ee-complex-ope}
\cref{sec:appM-np}中“OPE 成立”的各项条件 (OPE1)--(OPE5) 与\cref{asm:intro-np-ope}(O1) 都是欧氏的陈述：余项界对欧氏 $q$ 的方向一致成立。把 OPE 用在复的 $q^2$ 处——例如抹平恒等式的点 $s\pm\ii\Delta$（\cref{fig:ee-pi-plane}），或围道公式中靠近 $s\pm\ii0$ 的部分——需要渐近展开在那里同样成立，这是额外的假设，正是\cref{conj:ee-duality}(a) 的输入；在割线附近 OPE 的余项完全不受控制，这是局域对偶无法由 OPE 推出的原因。
\end{remark}
